\documentclass[11pt,a4paper]{article}
\usepackage[T1]{fontenc}
\usepackage{lmodern,amsmath,amssymb,amsthm,microtype}
\usepackage[margin=25mm]{geometry}
\usepackage[colorlinks=true,urlcolor=blue,linkcolor=blue]{hyperref}
\newtheorem{proposition}{Proposition}
\newcommand{\dd}{\mathop{}\!\mathrm d}
\title{A Barnes integral for the fixed-mass recoil graph}
\author{Alan Day \and Chengheng\thanks{AI research assistant (OpenAI).}}
\date{Working research note --- 30 October 2026}
\begin{document}\maketitle
\begin{abstract}
We reduce the full internal double-sine integral of the continued
tetrahedron graph to a Barnes integral, at fixed positive edge masses
in the pure-recoil limit. The transported contour becomes a horizontal
line separating two pole families. Exact \(q\)-Gamma factorization and
a uniform integrable tail bound justify passage of the limit through
this integral. Its vertical-contour representation yields four
convergent \({}_4F_3(1)\) terms for distinct pole families, while the
integral remains regular when they coalesce. Combining this result
with the previous vertex limit gives normalized fixed-topology graph
ratios. Physical external-energy cycles, insertion normalization and
the joint heavy-mass saddle limit remain uncomputed. This is an
author self-reviewed draft, without independent certification or a
claim of priority.
\end{abstract}

\section{Labels, contour and the exact change of variable}
Take the de Sitter radius to be one and use CLP~\cite{CLP}, (3.46).
The six edge masses are \(\mu_j>0\). Put
\[
 b^2=i/\kappa,\quad L=b^{-2}=-i\kappa,\quad Q=b+b^{-1},
 \qquad iP_j=-\frac1{2b}-\frac{ib\mu_j}{2}.
\]
Let
\[
 \mathcal V=\{123,156,246,345\},\qquad
 \mathcal R=\{1245,1346,2356\},\qquad
 a_v=\frac12\sum_{j\in v}\mu_j,\quad
 r_\rho=\frac12\sum_{j\in\rho}\mu_j .
\]
Each edge occurs twice in both collections, so \(\sum_v a_v=\sum_\rho r_\rho\).
Assume the strict separation
\begin{equation}
 a_{\max}<r_{\min},\qquad c=(a_{\max}+r_{\min})/2. \label{eq:gap}
\end{equation}
All masses and this gap are fixed as \(\kappa\to\infty\).

In the internal integral set
\begin{equation}
 iP=-\frac3{2b}+bz,\qquad \dd P=-ib\,\dd z. \label{eq:var}
\end{equation}
Its exact integrand, before the measure, is
\begin{equation}
 K_\kappa(z)=
 \frac{\prod_v S_b(b(z+ia_v))}
 {\prod_\rho S_b(b(z+\frac12+ir_\rho+L))\,
  S_b(b(z+\frac12-L))}. \label{eq:K}
\end{equation}
The original numerator pole families and denominator-zero families
become, with \(n,m\ge0\),
\begin{align}
 z&=-n+i\kappa m-ia_v,\\
 z&=\tfrac12+n-i\kappa m-ir_\rho,\qquad
 z=\tfrac12+n-i\kappa(m+2). \label{eq:poles}
\end{align}
Cancellations may remove poles; these lists contain all possible poles.

CLP (3.58), at the endpoint, gives
\(s=\mathbb R+\pi/\kappa-8\pi i+2\pi i(a_{\max}+r_{\min})/\kappa\).
Using \(s=-8\pi i-4\pi(z-\frac12)/\kappa\), its image is exactly
\begin{equation}
 \mathcal C_z=\mathbb R-ic. \label{eq:horizontal}
\end{equation}
For sufficiently large \(\kappa\), \eqref{eq:gap} puts the first
families above and the latter families below this line.
We orient \(\mathcal C_z\) left to right and use the corresponding
transported orientation in \eqref{eq:var}; reversing both changes only
an overall sign. This is an internal contour, not an external
three-energy contour for a physical correlator.

\section{Exact \texorpdfstring{\(q\)}{q}-Gamma factorization}
Write \(q=e^{-2\pi/\kappa}\), \(\tau=e^{-2\pi\kappa}\), and define
\[
 A_v=z+ia_v,\quad D_\rho=z+\tfrac12+ir_\rho,\quad D_0=z+\tfrac12,
 \qquad
 \Delta_2=\sum_v(A_v^2-A_v)-\sum_{\rho,0}(D_\rho^2-D_\rho).
\]
The last sum includes all four \(D\)'s. Then
\[
 \sum_v A_v-\sum_{\rho,0}D_\rho=-2,\qquad
 \Delta_2=-4z+\sum_\rho r_\rho^2-\sum_v a_v^2+1-i\sum_\rho r_\rho .
\]
CLP (A.5) and the \(q\)-Gamma definition~\cite{DLMFq} give exactly
\begin{equation}
 S_b(bw)=C_\kappa(1-q)^{w-1}
 e^{i\pi w/2+\pi(w^2-w)/(2\kappa)}
 \Gamma_q(w)(\tau e^{2\pi iw};\tau)_\infty, \label{eq:sb}
\end{equation}
where \(C_\kappa=(q;q)_\infty^{-1}e^{-i\pi(Q^2+1)/12}\)
is independent of \(w\).
The inverse-\(b\) shift is
\(S_b(bw+b^{-1})=2\sin(\pi w)S_b(bw)\).
Apply it once to each of the three shifted denominators and inversely
to the fourth. The normalized exact integrand is
\begin{align}
 \mathcal B_\kappa(z)&=-8\pi^3 e^{-\pi\kappa}(1-q)^2K_\kappa(z),\\
 &=\frac{\pi^3e^{-i\pi z+\pi\Delta_2/(2\kappa)}
                  (1+\tau e^{2\pi iz})}
              {\prod_\rho\sin(\pi D_\rho)}
 \frac{\prod_v\Gamma_q(A_v)}{\prod_{\rho,0}\Gamma_q(D_\rho)}
 \frac{\prod_v(\tau e^{2\pi iA_v};\tau)_\infty}
      {\prod_{\rho,0}(\tau e^{2\pi iD_\rho};\tau)_\infty}.
 \label{eq:Bq}
\end{align}
No small-\(b\) approximation has been made in this formula.
On compact subsets avoiding poles, the ordinary \(q\to1^-\) limit,
reflection formula and the small-\(\tau\) products yield
\begin{equation}
 \mathcal B_\kappa(z)\longrightarrow
 \mathcal B(z):=e^{-i\pi z}
 \frac{\prod_v\Gamma(z+ia_v)
       \prod_\rho\Gamma(\frac12-z-ir_\rho)}
      {\Gamma(z+\frac12)}. \label{eq:B}
\end{equation}

\section{An integrable bound for both tails}
Compact convergence alone would not establish the integral limit.
Here the mass balance supplies the missing tail control.

For \(q=e^{-h}\), put \(f_q(x)=\log\Gamma_q(x)\) at real \(x>0\).
Differentiating the absolutely convergent product and comparing its
positive sum with an integral gives
\begin{align}
 \log\frac{1-q^x}{1-q}-\frac{h}{e^{hx}-1}
 &\le f_q'(x)\le\log\frac{1-q^x}{1-q},\\
 |f_q''(x+s)|
 &\le \frac{h}{e^{h(x-A)}-1}
       +\frac{h^2e^{h(x-A)}}{(e^{h(x-A)}-1)^2}
 \le\frac1{x-A}+\frac1{(x-A)^2},
 \label{eq:qbounds}
\end{align}
provided \(\operatorname{Re}s\ge-A\) and \(x>A\). The complex bound
uses \(|1-q^{x+s+n}|\ge1-q^{x-A+n}\).

Thus, for any two fixed collections of four shifts with difference of
sums \(-2\), Taylor expansion along the shifts gives, uniformly for
large \(x\) and \(0<h\le h_0\),
\begin{equation}
 \left|\frac{\prod_{j=1}^4\Gamma_q(x+\alpha_j)}
                  {\prod_{j=1}^4\Gamma_q(x+\delta_j)}\right|
 \le C\left(\frac{1-q}{1-e^{-hx}}\right)^2. \label{eq:qratio}
\end{equation}
The constant depends on the shifts. In \eqref{eq:Bq}, on
\(z=x-ic\), all sine denominators stay away from zero, and the modular
products and their inverses are bounded for sufficiently large
\(\kappa\). The quadratic exponential contributes \(e^{-hx}\).
Consequently its positive tail is bounded by a constant times
\[
 e^{-hx}\left(\frac{1-e^{-h}}{1-e^{-hx}}\right)^2
 \le\left(\frac{h}{2\sinh(hx/2)}\right)^2\le x^{-2}.
\]

For clarity the negative tail is not inferred from the positive one.
Reflect all eight original double sines via \(S_b(Q-x)=1/S_b(x)\)
and put \(w=-z\). The exact integrand becomes
\begin{equation}
 K_\kappa(-w)=
 \frac{\prod_\rho S_b(b(w+\frac12-ir_\rho))\,
       S_b(b(w+\frac12+2L))}
      {\prod_v S_b(b(w+1-ia_v+L))}. \label{eq:negative}
\end{equation}
After inverse-\(b\) shifts its four small numerator arguments have
offsets \(\frac12-ir_\rho,\frac12\), and its four denominator
offsets are \(1-ia_v\). Their sum difference is again \(-2\).
The shift factors are
\[
 \frac{[2\sin\pi(w+\frac12+L)][2\sin\pi(w+\frac12)]}
      {\prod_v2\sin\pi(w+1-ia_v)}.
\]
On \(w=x+ic\), the denominator imaginary parts are \(c-a_v>0\).
Removing \(e^{\pi\kappa}\) leaves a bounded trigonometric factor,
and the same argument proves the negative-tail bound. Hence
\begin{equation}
 |\mathcal B_\kappa(x-ic)|\le\frac{C}{1+x^2}
 \quad(\kappa\ge\kappa_0,\ x\in\mathbb R). \label{eq:dom}
\end{equation}
On the remaining compact interval use \eqref{eq:Bq} and compact
convergence. The constants here depend on the fixed masses and gap;
no joint growing-mass bound is claimed.

\begin{proposition}
Under \eqref{eq:gap}, the normalized internal integral satisfies
\begin{equation}
 -\frac{4\pi^2(1-q)^2e^{-\pi\kappa}}{b}\,\mathcal J_b
 \longrightarrow
 \mathfrak B(\boldsymbol\mu):=
 \frac1{2\pi i}\int_{\mathbb R-ic}\mathcal B(z)\,\dd z
 =\frac1{2\pi i}\int_{d-i\infty}^{d+i\infty}\mathcal B(z)\,\dd z,
 \quad 0<d<\tfrac12. \label{eq:limit}
\end{equation}
\end{proposition}
\begin{proof}
Dominated convergence follows from \eqref{eq:dom}; the normalization
uses \(\dd P=-ib\,\dd z\) and the definition of \(\mathcal B_\kappa\).
For the second equality rotate the positive horizontal tail upwards
and the negative tail downwards about \(d-ic\). No poles enter either
swept quadrant. Gamma reflection and fixed-shift asymptotics give
\(O(|z|^{-2})\) on the connecting arcs, with exponential improvement
away from their horizontal ends, so the arc contributions vanish.
The vertical contour separates \(z=-n-ia_v\) from
\(z=\frac12+n-ir_\rho\). Its two tails decay as
\(e^{-2\pi y}\) for \(y\to+\infty\) and \(e^{-4\pi|y|}\) for
\(y\to-\infty\), times powers of \(|y|\).
\end{proof}

\section{A hypergeometric evaluation and coincident families}
For distinct \(a_v\), close the vertical contour to the left and
sum the four pole families. Standard Barnes residue calculus
\cite{DLMFb} gives
\begin{align}
 \mathfrak B&=\sum_{k=1}^4 c_k\,
 {}_4F_3\left(
 \begin{matrix}
 \frac12+ia_k,\ \{\frac12+i(a_k-r_\rho)\}_{\rho=1}^3\\
 \{1+i(a_k-a_j)\}_{j\ne k}
 \end{matrix};1\right), \label{eq:hyper}\\
 c_k&=e^{-\pi a_k}
 \frac{\prod_{j\ne k}\Gamma(i(a_j-a_k))
       \prod_\rho\Gamma(\frac12+i(a_k-r_\rho))}
      {\Gamma(\frac12-ia_k)}.
\end{align}
The sum of lower parameters minus the sum of upper parameters is one,
so each series converges. For example
\(\Gamma(A-n)=(-1)^n\Gamma(A)/(1-A)_n\);
all six signs from the residue, phase, other left factors and
reciprocal denominator cancel. Closing arcs can be chosen away from
the pole rows, where the balanced quotient is \(O(|z|^{-2})\).

Individual coefficients become singular when some \(a_j=a_k\).
The integral \eqref{eq:limit} is regular there and is the definition
to use, or one may take the limit of the complete sum.
Assigning values to the divergent summands separately is invalid.
Barnes/Racah representations have substantial antecedents; we make
no novelty claim or identification with a particular unitary group
\(6j\) without its representation and measure dictionary.

\section{What this now determines, and what it does not}
Let \(T(\mu_i,\mu_j,\mu_k)\) be the limiting combined vertex from
the preceding note. For a reference six-mass configuration
\(\boldsymbol\mu^\star\) satisfying \eqref{eq:gap} and
\(\mathfrak B(\boldsymbol\mu^\star)\ne0\), we obtain the full
fixed-topology graph ratio
\begin{equation}
 \lim_{\kappa\to\infty}
 \frac{Z_{\rm tet}(\boldsymbol\mu;\kappa)}
      {Z_{\rm tet}(\boldsymbol\mu^\star;\kappa)}
 =
 \prod_{v\in\mathcal V}\frac{T(\boldsymbol\mu_v)}
                            {T(\boldsymbol\mu^\star_v)}
 \left|\frac{\mathfrak B(\boldsymbol\mu)}
              {\mathfrak B(\boldsymbol\mu^\star)}\right|^2.
 \label{eq:graph}
\end{equation}
This is a graph result at fixed real masses, not a physical
four-point normalization prescription.

To connect to the prior transverse coordinates, choose explicitly
\[
 (\mu_1,\mu_2,\mu_3,\mu_4,\mu_5,\mu_6)
 =(m_1,\nu,m_2,m_3,\nu,m_4).
\]
Then the four \(a_v\)'s are the four values
\(M+\nu/2+(\pm u\pm v)/2\), independent of \(\omega\); the three
\(r_\rho\)'s are \(M+\nu+\omega/2,\,2M,\,M+\nu-\omega/2\).
At \(u=v=0\) all four left pole families coalesce, so
\eqref{eq:hyper} is not a stable termwise formula.
This explicit labeling may reverse the sign of \(\omega\) relative
to CLP's symmetric diagram; the earlier Gram expressions are even
in that sign.

The next calculable step is the large-\(M,\nu\) saddle analysis of
\eqref{eq:limit}, retaining \(u,v\), with the full vertex factor.
Its internal determinant must be derived in this representation
before comparison with the previous full-\(6j\) Gram hypothesis.
The tail constant in \eqref{eq:dom} has not been made uniform in
growing masses or complex external energies. Analytic continuation
of \eqref{eq:graph} requires separate chiral continuation, not a new
modulus at complex masses. The physical insertion factors, external
energy cycles and disconnected baseline remain undetermined.
Finite-\(G\) dynamics and a complete de Sitter dual remain open.

\section{Numerical checks}
The script \texttt{code/recoil\_internal.py} records 34 passing checks:
four symbolic balances, nine exact shift identities and precision
refinements, three limiting integrals at 45/65 digits, one
hypergeometric comparison, and eight two-sided tail refinements.
The four-term sum has cancellation ratio \(2.68\times10^5\) in its
sample; coalesced families are evaluated by the integral.
Relative convergence errors are saved separately from arithmetic
agreement. These are not interval certificates or direct full
finite-\(\kappa\) quadratures. The uniform tail result uses the
analytic argument above.

\begin{thebibliography}{9}
\bibitem{CLP} J. Chandra, D. Liska and B. Post,
\emph{3D de Sitter from complex Virasoro TQFT: Observers and knots},
\href{https://arxiv.org/html/2610.30247v1}{arXiv:2610.30247v1},
(3.46), (3.48), (3.52), (3.58), Appendix A.
\bibitem{DLMFq} NIST DLMF,
\href{https://dlmf.nist.gov/5.18}{Section 5.18},
\(q\)-Gamma product and classical limit; read online.
\bibitem{DLMFb} NIST DLMF,
\href{https://dlmf.nist.gov/16.5}{Section 16.5},
Barnes representations and residue continuation; read online.
\end{thebibliography}
\end{document}

